\begin{theorem}[Reconstruction from regular inverse coordinates]%
    \label{thm:peps_regular_parent_reconstruction}
    \lean{TNLean.PEPS.globalRegularTensorMap,
        TNLean.PEPS.regularSiteRetraction_eq_self_of_mem_range,
        TNLean.PEPS.globalRegularTensorMap_inverse_eq_self_of_mem_vertexImage,
        TNLean.PEPS.globalRegularTensorMap_inverse_eq_self_of_mem_parent,
        TNLean.PEPS.globalDependentPhysicalMap_regularInverse_leftMul}
    \leanok
    \uses{thm:peps_region_vertex_image_support,
        def:peps_vertex_image_ground_space}
    Let $A_v$ be regular $G$-injective site maps on a finite graph,
    and let $F_vA_v=\Pi_v$, where $\Pi_v$ averages simultaneous
    left translation of the incident group labels. Set
    $A_{\mathrm{all}}=\bigotimes_v A_v$ and
    $F_{\mathrm{all}}=\bigotimes_v F_v$.
    If a collection of parent regions covers every vertex and
    $\Psi$ satisfies each regional parent condition, then
    \begin{align}
        A_{\mathrm{all}}F_{\mathrm{all}}\Psi & =\Psi,
        \label{eq:peps_regular_parent_reconstruction}                           \\
        (F_{\mathrm{all}}\Psi)(k\cdot\alpha)
                                             & =(F_{\mathrm{all}}\Psi)(\alpha).
        \label{eq:peps_regular_parent_left_invariance}
    \end{align}
    Here $k\in G^V$ and $(k\cdot\alpha)_{v,e}=k_v\alpha_{v,e}$.
    This uses the restricted inverse in
    \cite[Definition~5.1]{Schuch2010PEPS}, with no injectivity
    requirement outside the invariant virtual subspace.
\end{theorem}
\begin{proof}\leanok
    Every one-site slice of $\Psi$ belongs to $\operatorname{im}A_v$.
    Since $A_v\Pi_v=A_v$, the composition $A_vF_v$ fixes that image.
    In the following local contraction, one virtual wire carries the
    full tuple of incident labels at $v$; the physical wire carries one
    site index. The inverse is used only after $A_v$:
    % Formula: A_v F_v A_v = A_v Pi_v = A_v, not A_v F_v = I globally.
    % Source: SCP10 Definition 5.1 and paper_v3.tex:1480--1514 (closure inverse).
    % Ink-to-index: matrices read right to left; each virtual wire bundles
    % eta_v in G^{Inc(v)}. A has physical row and virtual column; F the reverse.
    % Joined physical/virtual indices are summed. The external physical row
    % and virtual column remain open: every panel has (V,P)=(1,1).

    \begin{tenkzequation}
        \tenkzkernel
        \begin{tenkz}[rows={wire}, cols=3, bonds=none]
            \tn[at=(1,1), name=parentInverse0,
                ports={180:physical:$i$,0:virtual}]{$A_v$}
            \tn[at=(1,2), name=parentInverse1,
                ports={180:virtual,0:physical}]{$F_v$}
            \tn[at=(1,3), name=parentInverse2,
                ports={180:physical,0:virtual:$\eta_v$}]{$A_v$}
            \tnwire{parentInverse0.0}{parentInverse1.180}
            \tnwire{parentInverse1.0}{parentInverse2.180}
        \end{tenkz}
        $=$
        \begin{tenkz}[rows={wire}, cols=2, bonds=none]
            \tn[at=(1,1), name=parentAverage0,
                ports={180:physical:$i$,0:virtual}]{$A_v$}
            \tn[at=(1,2), name=parentAverage1,
                ports={180:virtual,0:virtual:$\eta_v$}]{$\Pi_v$}
            \tnwire{parentAverage0.0}{parentAverage1.180}
        \end{tenkz}
        $=$
        \begin{tenkz}[rows={wire}, cols=1, bonds=none]
            \tn[at=(1,1), name=parentRetracted0,
                ports={180:physical:$i$,0:virtual:$\eta_v$}]{$A_v$}
        \end{tenkz}
    \end{tenkzequation}
    The product of these independent retractions therefore fixes
    $\Psi$, proving~(\ref{eq:peps_regular_parent_reconstruction}).
    For $\phi=F_{\mathrm{all}}\Psi$, applying $F_{\mathrm{all}}$
    to~(\ref{eq:peps_regular_parent_reconstruction}) gives
    $\phi=(\bigotimes_v\Pi_v)\phi$. The rows of each $\Pi_v$
    are invariant under left translation, which gives
    (\ref{eq:peps_regular_parent_left_invariance}).
\end{proof}

\begin{theorem}[Shared-edge invariance and flat support of parent coordinates]%
    \label{thm:peps_regular_parent_coordinate_support}
    \lean{TNLean.PEPS.dependentRegionSlice_regularInverse_halfEdgeRightMul,
        TNLean.PEPS.dependentRegionSlice_regularInverse_apply_eq_zero_of_cycle_ne_one,
        TNLean.PEPS.globalDependentPhysicalMap_regularInverse_singleEdgeRightMul,
        TNLean.PEPS.globalDependentPhysicalMap_regularInverse_rightMul,
        TNLean.PEPS.globalDependentPhysicalMap_regularInverse_eq_zero_of_holonomy_ne_one,
        TNLean.PEPS.globalDependentPhysicalMap_regularInverse_eq_zero_of_torusPlaquetteHolonomy,
        TNLean.PEPS.mem_torusPlaquetteRegion_self,
        TNLean.PEPS.exists_torusPlaquetteRegion_contains_edge,
        TNLean.PEPS.IsGInjective.exists_torusParentCoordinates}
    \leanok
    \uses{thm:peps_regular_parent_reconstruction,
        lem:peps_regular_region_bond_right_invariance,
        thm:peps_regular_region_parent_flatness}
    In the preceding setting, suppose each edge has both endpoints in
    some parent region. Then $\phi=F_{\mathrm{all}}\Psi$ satisfies
    $\phi(\alpha\cdot r)=\phi(\alpha)$ for every $r\in G^E$, where
    $(\alpha\cdot r)_{v,e}=\alpha_{v,e}r_e$.
    Orient each edge from its smaller to its larger endpoint and define
    $u_e(\alpha)=\alpha_{\mathrm{head}(e),e}
        \alpha_{\mathrm{tail}(e),e}^{-1}$.
    If a closed walk $p$ lies in a connected parent region and its
    reverse-ordered holonomy satisfies $\operatorname{Hol}_{u(\alpha)}(p)\ne1$,
    then $\phi(\alpha)=0$.

    Every native vertex belongs to its own plaquette, and every native
    edge has both endpoints in some plaquette.
    In particular, for the native four-site plaquette parents on a
    regular $G$-injective torus with periods at least three, the maps
    $F_v$ can be chosen so that reconstruction, both translation
    invariances, and support on identity plaquette holonomies hold
    for every physical parent ground vector.
\end{theorem}
\begin{proof}\leanok
    Fix the exposed labels outside a region $R$. Its slice is a
    linear combination of the original physical slices after applying
    $\bigotimes_{v\in R}F_v$; the coefficients are products of the
    exterior inverse coefficients. Each transformed physical slice
    has shared internal-edge right invariance and vanishes at
    nonidentity cycle coordinates. Both properties pass to the
    linear combination. A single-edge right translation can therefore
    be performed in any parent region containing its endpoints.
    Successively translating the finitely many edges gives arbitrary $r$.

    If all cycle coordinates of a regional label are the identity,
    its internal quotient has the form $u_e=q_{\mathrm{head}(e)}
        q_{\mathrm{tail}(e)}^{-1}$. Its closed-walk holonomies telescope to
    $1$. Hence nonidentity holonomy implies a nonidentity cycle
    coordinate and forces the coefficient to vanish. Native plaquettes
    cover every vertex and every edge, and each plaquette is connected.
    The local inverses exist by regular $G$-injectivity.
\end{proof}

\begin{theorem}[Global spanning of the regular torus parent kernel]%
    \label{thm:peps_regular_torus_parent_spanning}
    \lean{TNLean.PEPS.globalRegularTensorMap_regularProjectorClosedState,
        TNLean.PEPS.IsGInjective.stateCoeff_mem_commutingClosureSpan_of_flat,
        TNLean.PEPS.IsGInjective.torusParentKernel_le_commutingClosureSpan}
    \leanok
    \uses{thm:peps_regular_parent_coordinate_support,
        thm:peps_regular_invariant_state_span,
        thm:peps_regular_closed_physical_contraction,
        thm:peps_flat_insertion_closure_state}
    Let $A$ be regular $G$-injective on a native torus whose two
    periods are at least three. For each native plaquette $R_v$,
    let $h_v\ge0$ have kernel equal to the open-region ground space
    on $R_v$, and let $H=\sum_v h_v\otimes I_{R_v^c}$. Then
    \begin{align}
        \ker H\subseteq
        \operatorname{span}\{\lvert\Psi_{g,h}\rangle:gh=hg\}.
        \label{eq:peps_regular_global_spanning}
    \end{align}
    This is the global spanning direction of
    \cite[Theorem~5.7]{Schuch2010PEPS} for regular virtual
    representations. No membership in a closure span is assumed.
\end{theorem}
\begin{proof}\leanok
    Positivity of the local terms implies that each annihilates any
    $\Psi\in\ker H$. Choose inverse coordinates $\phi=F_{\mathrm{all}}\Psi$.
    Let $\mathcal F$ be the set of edge assignments with identity
    holonomy on every native plaquette, let $\alpha^u$ have tail label
    $1$ and head label $u_e$, and let $C_u$ be the closed contraction
    of regular averaging projectors with insertion $u$.
    The supported orbit expansion and reconstruction give
    \begin{align}
        \Psi=A_{\mathrm{all}}\phi
         & =\sum_{u\in\mathcal F}\phi(\alpha^u)A_{\mathrm{all}}C_u
        =\sum_{u\in\mathcal F}\phi(\alpha^u)\lvert\Psi_u\rangle.
        \label{eq:peps_regular_flat_expansion}
    \end{align}
    Put $\Omega_u=\sum_{\eta\in G^E}\lvert\alpha(u,\eta)\rangle$,
    where $\alpha(u,\eta)$ has tail label $\eta_e$ and head label
    $u_e\eta_e$. This bond vector is unnormalized. With
    $\Pi_{\mathrm{all}}=\bigotimes_v\Pi_v$, one has
    $C_u=\Pi_{\mathrm{all}}\Omega_u$ and
    $A_{\mathrm{all}}\Pi_{\mathrm{all}}=A_{\mathrm{all}}$.
    Each $\Pi_v$ includes the averaging factor $|G|^{-1}$, so
    $\Pi_{\mathrm{all}}$ includes $|G|^{-|V|}$.
    Thus the projector factors in (\ref{eq:peps_regular_flat_expansion})
    are absorbed site by site:
    % Formula: sum_{u flat} phi(alpha^u) A_all Pi_all Omega_u
    %        = sum_{u flat} phi(alpha^u) A_all Omega_u.
    % Source: SCP10 Theorem 5.7; regular-coordinate proof in
    % TorusRegularParentSpanning.lean, via RegularInvariantStateSpan.lean.
    % Ink-to-index: virtual wire bundles ALL half-edge labels, not one bond;
    % physical wire bundles sigma in the full physical configuration space.
    % Pi_all includes |G|^{-|V|} sum_{k in G^V} L_k; Omega has no scalar factor.
    % All virtual labels contract. Each panel has (V,P)=(0,1), the one
    % physical port denoting the full tuple sigma, not a one-site output.

    \begin{tenkzequation}
        \tenkzkernel
        $\displaystyle\sum_{u\in\mathcal F}\phi(\alpha^u)$
        \begin{tenkz}[rows={wire}, cols=3, bonds=none]
            \tn[at=(1,1), name=parentFlat0,
                ports={180:physical:$\sigma$,0:virtual}]{$A_{\mathrm{all}}$}
            \tn[at=(1,2), name=parentFlat1,
                ports={180:virtual,0:virtual}]{$\Pi_{\mathrm{all}}$}
            \tn[at=(1,3), name=parentFlat2,
                ports={180:virtual}]{$\Omega_u$}
            \tnwire{parentFlat0.0}{parentFlat1.180}
            \tnwire{parentFlat1.0}{parentFlat2.180}
        \end{tenkz}
        $=$
        $\displaystyle\sum_{u\in\mathcal F}\phi(\alpha^u)$
        \begin{tenkz}[rows={wire}, cols=2, bonds=none]
            \tn[at=(1,1), name=parentAbsorbed0,
                ports={180:physical:$\sigma$,0:virtual}]{$A_{\mathrm{all}}$}
            \tn[at=(1,2), name=parentAbsorbed1,
                ports={180:virtual}]{$\Omega_u$}
            \tnwire{parentAbsorbed0.0}{parentAbsorbed1.180}
        \end{tenkz}
    \end{tenkzequation}
    Every $u\in\mathcal F$ is gauge equivalent to two commuting
    closure seams. Local regular invariance leaves the physical
    contraction unchanged under this vertex gauge. Thus every
    summand in~(\ref{eq:peps_regular_flat_expansion}) belongs to the
    right-hand side of~(\ref{eq:peps_regular_global_spanning}).
\end{proof}
